% NeuroVascNet progress slides: Muntasir and Maisha
% Four slides each, in order: done so far, found, current challenges, novelty.
% Compile with pdflatex (Overleaf works out of the box).
\documentclass[aspectratio=169,11pt]{beamer}

\usepackage[utf8]{inputenc}
\usepackage[T1]{fontenc}
\usepackage{lmodern}
\usepackage[most]{tcolorbox}
\usepackage{booktabs}

% ---------------------------------------------------------------- look
\usetheme{default}
\setbeamertemplate{navigation symbols}{}
\setbeamertemplate{footline}{%
  \hfill\scriptsize\insertframenumber\,/\,\inserttotalframenumber\hspace{1.2em}\vspace{0.8em}}
\definecolor{accent}{HTML}{2A78D6}
\definecolor{ink}{HTML}{1A1A1A}
\setbeamercolor{frametitle}{fg=ink,bg=white}
\setbeamerfont{frametitle}{size=\Large}
\setbeamercolor{normal text}{fg=ink}
\setbeamercolor{itemize item}{fg=accent}
\setbeamertemplate{itemize items}[circle]
\setlength{\leftmargini}{1.2em}

\newtcolorbox{slidebox}{enhanced, colback=white, colframe=accent!80!black,
  boxrule=0.9pt, arc=1pt, left=8pt, right=8pt, top=6pt, bottom=6pt, width=\textwidth}
\newcommand{\colhead}[1]{\textbf{#1}\par\vspace{2pt}}

\title{NeuroVascNet: Where We Are}
\subtitle{Can EEG tell us what the blood flow in the brain is doing?}
\author{Muntasir Hossain \and Maisha}
\date{October 2026}

\begin{document}

\begin{frame}
  \titlepage
\end{frame}

% ================================================================ MUNTASIR
\begin{frame}{Muntasir Experiment}
\begin{slidebox}
\begin{itemize}\setlength{\itemsep}{5pt}
  \item Before building any model, I checked the data itself: 610 recordings across our 4 datasets.
  \item That check showed HEFMI's fNIRS was mostly heartbeat and breathing, so I switched training to the cleaner TU Berlin dataset.
  \item I built one EEG model that learns from fNIRS and fMRI together, and a second one that learns from fNIRS alone, to compare them fairly.
  \item Every test uses people the model has never seen, with chance-level controls.
  \item In HEFMI, I compared 20 stroke (ICH) patients with 17 healthy people: how strongly and how quickly EEG and blood signals are linked, plus four checks for recording problems.
\end{itemize}
\end{slidebox}
\end{frame}

\begin{frame}{Muntasir Found}
\begin{slidebox}
\begin{columns}[T,onlytextwidth]
\begin{column}{0.48\textwidth}
\colhead{What worked}
\begin{itemize}\setlength{\itemsep}{4pt}
  \item EEG can predict fNIRS in new people, but only a little (r = 0.16), and only with clean data.
  \item The delay between brain activity and blood flow is about 6 seconds, and it stays the same for a person across sessions.
  \item Stroke patients' HbO and HbR move together more tightly than healthy people's (0.81 vs 0.64). The EEG--blood link itself was the same in both groups.
\end{itemize}
\end{column}
\begin{column}{0.48\textwidth}
\colhead{What didn't}
\begin{itemize}\setlength{\itemsep}{4pt}
  \item Adding fMRI did not help. Random fMRI data gave the same result as real fMRI.
  \item The model fails completely on a different dataset.
  \item The model cannot tell when its own prediction is wrong.
  \item Splitting the data the usual way made our score look almost twice as good as it really is.
\end{itemize}
\end{column}
\end{columns}
\end{slidebox}
\end{frame}

\begin{frame}{Muntasir Current Challenges}
\begin{slidebox}
\begin{itemize}\setlength{\itemsep}{5pt}
  \item No public dataset records the same people with EEG, fNIRS and fMRI, so the two blood signals never meet in the data.
  \item The fMRI dataset has only 10 people, which is very little to learn from.
  \item The outside datasets differ in electrodes, task and data quality all at once, so I can't tell which one breaks the model.
  \item The stroke result might come from how light passes through an injured head, not from the brain. I still need to rule that out.
  \item I haven't fully confirmed that EEG and blood recordings are perfectly lined up in time.
\end{itemize}
\end{slidebox}
\end{frame}

\begin{frame}{Muntasir Novelty Concern}
\begin{slidebox}
\begin{itemize}\setlength{\itemsep}{5pt}
  \item As far as I know, this is the first fair test of whether training with fMRI on top of fNIRS helps an EEG model. It doesn't.
  \item Checking data quality first predicts whether this kind of model can work at all: clean data gave r = 0.16, noisy data gave nothing.
  \item We show how much a common shortcut in testing inflates results.
  \item We measure the brain-to-blood delay directly, and show why the model's own delay estimate can't be trusted.
  \item A possible new marker for stroke patients, still to be confirmed.
\end{itemize}
\end{slidebox}
\end{frame}

% ================================================================ MAISHA
\begin{frame}{Maisha Experiment}
\begin{slidebox}
\begin{itemize}\setlength{\itemsep}{5pt}
  \item I rebuilt the pipeline independently, training on HEFMI fNIRS and ds002336 fMRI.
  \item I fixed the validation split: the old one had put only stroke patients into validation.
  \item I added safety checks: a test of whether global signal regression helps, unit and timing checks, and a rule that stops the notebook from claiming more than the data shows.
  \item I tested every person separately (leave one person out) and tried calibrating the model to each new person.
  \item I repeated the whole run 5 times with different random seeds to see how stable the results are.
\end{itemize}
\end{slidebox}
\end{frame}

\begin{frame}{Maisha Found}
\begin{slidebox}
\begin{columns}[T,onlytextwidth]
\begin{column}{0.48\textwidth}
\colhead{What worked}
\begin{itemize}\setlength{\itemsep}{4pt}
  \item EEG and fNIRS really are linked inside each person (r = 0.09, well above chance).
  \item Within our own dataset, the confidence score does put good predictions ahead of bad ones, in all 5 runs.
  \item Global signal regression throws away real signal, so I turned it off.
\end{itemize}
\end{column}
\begin{column}{0.48\textwidth}
\colhead{What didn't}
\begin{itemize}\setlength{\itemsep}{4pt}
  \item For a new person, the model does no better than guessing that person's usual level.
  \item Adding fMRI made things worse in 4 out of 5 runs.
  \item Nothing carried over to the other datasets.
  \item Personalising helped, but mostly because it learned the person's average, not because of EEG.
\end{itemize}
\end{column}
\end{columns}
\end{slidebox}
\end{frame}

\begin{frame}{Maisha Current Challenges}
\begin{slidebox}
\begin{itemize}\setlength{\itemsep}{5pt}
  \item The fNIRS data I train on is mostly body noise, which limits everything I can find.
  \item I only have 7 people to validate on (3 for fMRI), so results change from run to run.
  \item In half of the data splits, the timing between EEG and blood flow didn't look right, so I can't make timing claims yet.
  \item The "refuse to answer" rule works in one run and fails in the next.
  \item Each run takes about an hour and a half on our GPU, so testing ideas is slow.
\end{itemize}
\end{slidebox}
\end{frame}

\begin{frame}{Maisha Novelty Concern}
\begin{slidebox}
\begin{itemize}\setlength{\itemsep}{5pt}
  \item My pipeline is built separately from Muntasir's and reaches the same negative results, so they are not one pipeline's mistake.
  \item With 30 people, I can show the EEG--blood link is personal: it exists, but doesn't transfer between people.
  \item A simple check that shows how little EEG adds when personalising the model.
  \item A set of safeguards that keep the notebook honest about what it can claim.
  \item I found and fixed a split that would have mixed up patients and healthy people.
\end{itemize}
\end{slidebox}
\end{frame}

% ================================================================ DATASETS
\begin{frame}{The Data We Used}
\begin{slidebox}
\small
We both used the same four public datasets (about 85 GB in total), but not in the same roles. Before trusting any result, we checked each one.
\vspace{6pt}

\footnotesize
\begin{tabular}{@{}p{0.19\textwidth}p{0.29\textwidth}p{0.21\textwidth}p{0.23\textwidth}@{}}
\toprule
\textbf{Dataset} & \textbf{What it has} & \textbf{Muntasir used it for} & \textbf{Maisha used it for} \\
\midrule
TU Berlin hBCI & EEG + fNIRS, 29 people & training & testing (new dataset) \\
HEFMI-ICH & EEG + fNIRS, 37 people (20 stroke patients) & testing + stroke study & training and validation \\
ds002336 & EEG + fMRI, 10 people & training & training and validation \\
ds002739 & EEG + fMRI, 24 people & testing (new dataset) & testing (new dataset) \\
\bottomrule
\end{tabular}
\end{slidebox}
\end{frame}

% ---------------------------------------------------------------- TU Berlin
\begin{frame}{Dataset 1: TU Berlin hBCI (EEG + fNIRS) -- the data}
\begin{slidebox}
\scriptsize
\begin{columns}[T,onlytextwidth]
\begin{column}{0.485\textwidth}
\colhead{\small The recording}
\begin{tabular}{@{}p{0.3\linewidth}p{0.66\linewidth}@{}}
\toprule
Type & EEG and fNIRS recorded at the same time; a brain--computer interface study \\
Size & 7.1 GB \\
Participants & 29 healthy adults, no patients \\
Classes & 4: left vs right hand imagery; mental arithmetic vs rest \\
Trials & 3 sessions per task, 20 trials each: 120 per person, 3,480 in total \\
Conditions & seated at a screen; cue, about 10 s of task, then 15--17 s rest \\
\bottomrule
\end{tabular}
\end{column}
\begin{column}{0.485\textwidth}
\colhead{\small Signals and windows}
\begin{tabular}{@{}p{0.3\linewidth}p{0.66\linewidth}@{}}
\toprule
EEG & 30 electrodes (plus eye channels), 200 Hz, finer 10-5 positions (no C3/C4) \\
fNIRS & 36 channels, 2 wavelengths (760/850 nm), 10 Hz, raw light \\
Frequencies & EEG up to 45 Hz; fNIRS 0.01--0.2 Hz \\
Trial window & 30 s ($-$2 to +28 s around the cue): 6,000 EEG and 300 fNIRS samples \\
Alignment & both devices share the same event markers; windows cut at matching cue times \\
Model output & HbO and HbR: 72 signals $\times$ 300 time points \\
\bottomrule
\end{tabular}
\end{column}
\end{columns}
\end{slidebox}
\end{frame}

\begin{frame}{Dataset 1: TU Berlin hBCI -- what we found}
\begin{slidebox}
\small
\textbf{How we used it:} Muntasir trained on it; Maisha only tested on it.
\begin{itemize}\setlength{\itemsep}{3pt}
  \item The cleanest fNIRS data we have: it repeats well across trials (0.50) and behaves like a real brain response.
  \item Most optodes were touching the skin properly (about 78\% of channels).
  \item Its electrodes sit in different places from HEFMI's, so we had to estimate the missing ones.
  \item Its values were in different units from HEFMI's (thousands of times larger), so we converted them.
  \item Because it came as raw light, we did the conversion to HbO and HbR ourselves.
  \item A different 26-person TU Berlin dataset had been mixed up with this one early on; we caught it.
\end{itemize}
\end{slidebox}
\end{frame}

% ---------------------------------------------------------------- HEFMI
\begin{frame}{Dataset 2: HEFMI-ICH (EEG + fNIRS) -- the data}
\begin{slidebox}
\scriptsize
\begin{columns}[T,onlytextwidth]
\begin{column}{0.485\textwidth}
\colhead{\small The recording}
\begin{tabular}{@{}p{0.3\linewidth}p{0.66\linewidth}@{}}
\toprule
Type & EEG and fNIRS recorded at the same time; a clinical study \\
Size & 13.2 GB \\
Participants & 37: 17 healthy, 20 patients after a brain bleed (ICH) \\
Classes & 2: left vs right hand imagery \\
Trials & 3 sessions, 15 per hand per session: 90 per person, 2,940 used \\
Conditions & same protocol for patients and healthy people \\
\bottomrule
\end{tabular}
\end{column}
\begin{column}{0.485\textwidth}
\colhead{\small Signals and windows}
\begin{tabular}{@{}p{0.3\linewidth}p{0.66\linewidth}@{}}
\toprule
EEG & 32 electrodes, 256 Hz, standard 10-20 positions (C3, C4 present) \\
fNIRS & 90 channels, 11 Hz, already converted to HbO and HbR \\
Frequencies & EEG 0.5--45 Hz; fNIRS 0.01--0.2 Hz \\
Trial window & 30 s, resampled to 200 Hz EEG and 10 Hz fNIRS \\
Alignment & files come already cut into trials, EEG and fNIRS together, but without event times \\
Model output & HbO and HbR: 180 signals $\times$ 300 time points (Maisha also 4 summary numbers) \\
\bottomrule
\end{tabular}
\end{column}
\end{columns}
\end{slidebox}
\end{frame}

\begin{frame}{Dataset 2: HEFMI-ICH -- what we found}
\begin{slidebox}
\small
\textbf{How we used it:} Muntasir tested on it and compared patients with healthy people; Maisha trained on it (30 people) and validated on 7.
\begin{itemize}\setlength{\itemsep}{3pt}
  \item The fNIRS signal is mostly heartbeat and breathing, not brain activity, and it barely repeats across trials (0.08).
  \item Our early cleaning steps made it even worse.
  \item The original validation split had only patients in it; Maisha rebuilt it with both groups.
  \item The electrode names were missing, so we recovered them from the dataset's own code.
  \item With no event times, some checks were impossible.
  \item Timing looked fine for the training people but not for the validation people.
\end{itemize}
\end{slidebox}
\end{frame}

% ---------------------------------------------------------------- ds002336
\begin{frame}{Dataset 3: ds002336 (EEG + fMRI) -- the data}
\begin{slidebox}
\scriptsize
\begin{columns}[T,onlytextwidth]
\begin{column}{0.485\textwidth}
\colhead{\small The recording}
\begin{tabular}{@{}p{0.3\linewidth}p{0.66\linewidth}@{}}
\toprule
Type & EEG recorded inside an MRI scanner; a neurofeedback study \\
Size & 19.3 GB \\
Participants & 10 healthy adults, no patients \\
Classes & rest vs motor imagery blocks \\
Tasks & motor localiser, imagery before and after training, three neurofeedback runs \\
Trials & up to 100 per run; 384 used (Maisha 305 training + 124 validation) \\
\bottomrule
\end{tabular}
\end{column}
\begin{column}{0.485\textwidth}
\colhead{\small Signals and windows}
\begin{tabular}{@{}p{0.3\linewidth}p{0.66\linewidth}@{}}
\toprule
EEG & 64 electrodes at 5,000 Hz, cleaned of scanner noise, resampled to 200 Hz \\
fMRI & whole brain, one scan every 2 s (0.5 Hz), averaged into brain regions \\
Frequencies & EEG 0.5--45 Hz; fMRI 0.008--0.15 Hz \\
Trial window & 30 s: 6,000 EEG samples and 15 brain scans \\
Alignment & the scanner writes a marker into the EEG at every scan, so each scan matches an exact EEG sample \\
Model output & BOLD signal per brain region $\times$ 15 time points \\
\bottomrule
\end{tabular}
\end{column}
\end{columns}
\end{slidebox}
\end{frame}

\begin{frame}{Dataset 3: ds002336 -- what we found}
\begin{slidebox}
\small
\textbf{How we used it:} the only fMRI training set for both of us; Maisha also kept 3 people back for validation.
\begin{itemize}\setlength{\itemsep}{3pt}
  \item The raw EEG was buried under scanner noise about 500 times bigger than the brain signal; our own clean-up brought it back to normal levels.
  \item A few recordings had no matching partner file or were completely flat, so we dropped them (52 good pairs left).
  \item People barely moved their heads, which is good.
  \item The motor brain response is weak, and part of it is just a visual response to the cue.
  \item With only 10 people, most of what looked like prediction was the model recognising who the person was.
\end{itemize}
\end{slidebox}
\end{frame}

% ---------------------------------------------------------------- ds002739
\begin{frame}{Dataset 4: ds002739 (EEG + fMRI) -- the data}
\begin{slidebox}
\scriptsize
\begin{columns}[T,onlytextwidth]
\begin{column}{0.485\textwidth}
\colhead{\small The recording}
\begin{tabular}{@{}p{0.3\linewidth}p{0.66\linewidth}@{}}
\toprule
Type & EEG recorded inside an MRI scanner; a decision-making study \\
Size & 45.4 GB (our largest) \\
Participants & 24 healthy adults, aged 20--32, no patients \\
Classes & no motor classes: decide which way dots move, then rate confidence \\
Runs & 48 (two per person) \\
Trials & 4,678 used (Maisha 2,448) \\
\bottomrule
\end{tabular}
\end{column}
\begin{column}{0.485\textwidth}
\colhead{\small Signals and windows}
\begin{tabular}{@{}p{0.3\linewidth}p{0.66\linewidth}@{}}
\toprule
EEG & 64 electrodes recorded at 5,000 Hz; shared already cleaned at 1,000 Hz; no electrode names \\
fMRI & one scan every 2 s (0.5 Hz), 70 $\times$ 70 $\times$ 32 voxels, averaged into brain regions \\
Frequencies & EEG 0.5--40 Hz (filtered by the authors) \\
Trial window & 30 s: 15 brain scans per window \\
Alignment & event times from the dataset; EEG channels matched only by their order \\
Model output & BOLD signal per brain region $\times$ 15 time points \\
\bottomrule
\end{tabular}
\end{column}
\end{columns}
\end{slidebox}
\end{frame}

\begin{frame}{Dataset 4: ds002739 -- what we found}
\begin{slidebox}
\small
\textbf{How we used it:} an unseen fMRI test set for both of us, never used for training.
\begin{itemize}\setlength{\itemsep}{3pt}
  \item The best image quality of all our data.
  \item But it is a decision task, not a motor task, so it tests a different kind of brain activity from what we trained on.
  \item The quick decisions give only a small, inconsistent blood response.
  \item Without electrode names, we could only match channels by their order.
  \item One of our automatic quality rules wrongly flagged almost everyone, because the authors had already filtered the data. We reported it rather than quietly changing the rule.
\end{itemize}
\end{slidebox}
\end{frame}

\end{document}
